\section*{3.3 Lebesgue--Stieltjes Measures}

Given a Stieltjes measure function, we can apply the measure extension theorem to construct a measure on the Borel algebra $\mathcal{B}(\mathbb{R})$. Intuitively, the Lebesgue--Stieltjes measure $\mu$ assigns a length to each Borel set on the real line, where the length of an interval $(a, b]$ is given by the difference in the values of $F$ at the endpoints $a$ and $b$.

\begin{thmbox}{3.3}
Let $F$ be a Stieltjes measure function defined on $\mathbb{R}$. Then there exists a unique measure $\mu$ defined on Borel $\sigma$-algebra $\mathcal{B}(\mathbb{R})$ such that $\mu((a, b]) = F(b) - F(a)$. The measure $\mu$ is called the \textit{Lebesgue--Stieltjes measure} induced by $F$.
\end{thmbox}
This theorem is an application of the measure extension theorem (Theorem 3.1). The proof relies on a standard result from real analysis, which we first recall below.

\begin{defbox}{3.6}
A closed and bounded interval in $\mathbb{R}$ is called a \textit{compact} set in $\mathbb{R}$.
\end{defbox}

\begin{defbox}{3.7}
An \textit{open cover} of a set $A \subseteq \mathbb{R}$ is a collection of open sets $\{A_\alpha\}_{\alpha \in I}$, indexed by an arbitrary set $I$, such that $A$ is contained in the union of $A_\alpha$ over all $\alpha \in I$. A \textit{subcover} of $A$ is a sub-collection of $\{A_\alpha\}_{\alpha \in I}$ that is also a cover of $A$. If a cover consists of finitely many open sets, it is said to be a \textit{finite cover}.
\end{defbox}

To illustrate these definitions, consider the closed interval $[0, 1]$, which is a compact set. Let $\epsilon > 0$ be a fixed real number. We can cover $[0, 1]$ by infinitely many open intervals of the form $(a - \epsilon, a + \epsilon)$, where $a$ ranges over all real numbers in $[0, 1]$. However, there is a lot of redundancy in this open cover, and we can in fact cover $[0, 1]$ using only finitely many open sets from this collection. For instance, we can choose the open intervals $(k\epsilon - \epsilon, k\epsilon + \epsilon)$, where $k$ ranges over $0, 1, 2, \dots, \lfloor 1/\epsilon \rfloor$.

The Heine--Borel theorem states that the above example is a general phenomenon. We can always find a finite subcover for any compact set.

\begin{thmbox}{3.4 (Heine--Borel)}
Let $A$ be a subset of $\mathbb{R}$. Then $A$ is compact if and only if every open cover of $A$ has a finite subcover.
\end{thmbox}

In the proof of the main result Theorem 3.3 in this section, it is more convenient to show that a set function is $\sigma$-additive through $\sigma$-subadditivity.

\begin{defbox}{3.8}
\begin{enumerate}[label=(\roman*)]
    \item A set function $\mu$ on a field $\mathcal{F}_0$ is said to be \textit{finitely additive} if for any finite collection of disjoint set $A_1, A_2, \dots, A_n$ in $\mathcal{F}_0$, we have
    \[
    \mu\Bigl(\biguplus_{i=1}^{n} A_i\Bigr) = \sum_{i=1}^{n} \mu(A_i).
    \]
    \item A set function $\mu$ on a field $\mathcal{F}_0$ is called \textit{$\sigma$-subadditive} if for any countable collection $(A_i)_{i \ge 1}$ of sets in $\mathcal{F}_0$ such that $\cup_{i=1}^{\infty} A_i \in \mathcal{F}_0$, we have
        \[
    \mu\Bigl(\bigcup_{i=1}^{\infty} A_i\Bigr) \le \sum_{i=1}^{\infty} \mu(A_i).
    \]
\end{enumerate}
\end{defbox}

One can show that any $\sigma$-additive function satisfies the two conditions in the above definition.

\begin{thmbox}{3.5}
A set function $\mu$ defined on a field $\mathcal{F}_0$, with $\mu(\emptyset) = 0$, is $\sigma$-additive if:
\begin{enumerate}
    \item $\mu$ is finitely additive, and
    \item $\mu$ is $\sigma$-subadditive.
\end{enumerate}
Hence, a nonnegative set function on a field $\mathcal{F}_0$ that satisfies $\mu(\emptyset) = 0$, finite additivity, and $\sigma$-subadditivity is a pre-measure on $\mathcal{F}_0$.
\end{thmbox}

\noindent\textit{Proof} Suppose $\mu: \mathcal{F}_0 \to [0, \infty]$ is a set function that satisfies finite additivity, $\sigma$-subadditivity, and $\mu(\emptyset) = 0$. Suppose $E_i$ are mutually disjoint sets in $\mathcal{F}_0$, for $i = 1, 2, 3, \dots$, such that $\uplus_{i=1}^{\infty} E_i = E \in \mathcal{F}_0$. For each finite $n$, we have
\[
\mu(E) \ge \mu(E_1 \uplus E_2 \uplus \dots \uplus E_n) = \sum_{i=1}^{n} \mu(E_i).
\]
Taking the limit as $n \to \infty$, we obtain $\mu(E) \ge \sum_{i=1}^{\infty} \mu(E_i).

To show the reverse inequality, we note that $\mu$ is $\sigma$-subadditivity, and so
\[
\mu(E) = \mu\Bigl(\bigcup_{i=1}^{\infty} E_i\Bigr) \le \sum_{i=1}^{\infty} \mu(E_i).
\]
This shows that $\mu$ is $\sigma$-additive and hence is a pre-measure. \hfill $\blacksquare$

\noindent\textit{Proof of Theorem 3.3} Suppose $F$ is a Stieltjes measure function and $\mathcal{B}_0$ is the collection of sets as defined in Example 3.1.2, which are disjoint unions of intervals in the form $(a, b]$. We define a set function on $\mathcal{B}_0$ as follows: for any $(a, b]$ in $\mathcal{B}_0$, we set $\mu_0((a, b]) \triangleq F(b) - F(a)$, and we define $\mu_0(\emptyset) = 0$. For any disjoint union $A = \uplus_{i=1}^{m} (a_i, b_i]$, we set
\[
\mu_0(A) \triangleq \sum_{i=1}^{m} [F(b_i) - F(a_i)].
\]
\textbf{Step 1.} Show that $\mu_0$ is a well-defined set function.

Suppose a set $A \in \mathcal{B}_0$ can be written in two ways:
\[
A = \uplus_{i=1}^m (a_i, b_i] = \uplus_{j=1}^n (c_j, d_j].
\]
We want to check that the value of $\mu_0$ on $A$ is independent of the choice of representation. The proof idea is to consider a common refinement of the two partitions.

For each pair $(i, j)$, for $1 \le i \le m$ and $1 \le j \le n$, define the interval
\[
I_{ij} = (a_i, b_i] \cap (c_j, d_j] = (\max(a_i, c_j), \min(b_i, d_j)],
\]
which may be empty. It can be verified that the intervals $I_{ij}$ are pair-wise disjoint and their union is equal to $A$. Furthermore, $\{I_{ij}\}_{i,j}$ is a subdivision of the partitions $\{(a_i, b_i]\}_{i=1}^m$ and $\{(c_j, d_j]\}_{j=1}^n$. Hence,
\[
\sum_{i=1}^n \mu_0((a_i, b_i]) = \sum_{i=1}^m \sum_{j=1}^n \mu_0(I_{ij}) = \sum_{j=1}^n \mu_0((c_j, d_j]).
\]
This shows that $\mu_0$ is well-defined on $\mathcal{B}_0$.

\textbf{Step 2.} Show that $\mu_0$ is finitely additive.

Since a set in the algebra $\mathcal{B}_0$ is a finite disjoint union of left-open-right-closed intervals, it is sufficient to consider the union of disjoint intervals $(a_i, b_i]$, for $i = 1, 2, \dots, m$ such that $\uplus_{i=1}^m (a_i, b_i]$ is equal to $(a, b]$. After some re-labeling, we can re-order them and write
\[
a = a_1 < b_1 = a_2 < b_2 = a_3 < b_3 \cdots a_n < b_n = b.
\]
We have a telescoping sum
\[
\sum_{i=1}^n \mu_0((a_i, b_i]) = \sum_{i=1}^n [F(b_i) - F(a_i)] = F(b) - F(a) = \mu_0((a, b]).
\]
This proves that $\mu_0$ is finitely additive.

\textbf{Step 3.} Show that $\mu_0$ is $\sigma$-subadditive.

Suppose $(a, b] = \bigcup_{i=1}^\infty (a_i, b_i]$. The subintervals $(a_i, b_i]$'s need not be mutually disjoint. We want to show
\begin{equation}
\mu_0((a, b]) \le \sum_{i=1}^\infty \mu_0((a_i, b_i]). \tag{3.2}
\end{equation}
The difficulty in proving (3.2) is that it involves infinitely many intervals. We would like to reduce it to finitely many intervals using the Heine--Borel theorem.


As in a typical proof of real analysis, we fix a positive but arbitrarily small $\epsilon$ and make some relaxation. Pick $\delta > 0$ such that $F(a + \delta) < F(a) + \epsilon$. We choose $\delta$ small enough so that $a + \delta < b$. Note that the right continuity of $F$ is used when we pick $\delta$.

For each $i = 1, 2, 3, \dots$, pick $\eta_i > 0$ such that $F(b_i + \eta_i) < F(b_i) + \frac{\epsilon}{2^i}$. We have an open cover
\[
[a + \delta, b] \subseteq \bigcup_{i=1}^\infty (a_i, b_i + \eta_i).
\]
By Heine--Borel theorem, there exist $i_1, i_2, \dots, i_m$, such that $(a_{i_j}, b_{i_j} + \eta_{i_j})$, for $j = 1, 2, \dots, m$, form a finite cover of the closed and bounded interval $[a + \delta, b]$,
\[
[a + \delta, b] \subseteq \bigcup_{j=1}^m (a_{i_j}, b_{i_j} + \eta_{i_j}).
\]
We cannot directly apply $\mu_0$ to both sides in the above line because $\mu_0$ is only defined on left-open-right-closed intervals, and their finite unions. To this end, we remove one point on the left side and add $m$ points on the right side to obtain
\[
(a + \delta, b] \subseteq \bigcup_{j=1}^m (a_{i_j}, b_{i_j} + \eta_{i_j}].
\]
We now apply the set function $\mu_0$ to both sides and use the assumption that $\mu_0$ is finitely additive. This gives
\[
F(b) - F(a + \delta) \le \sum_{j=1}^m F(b_{i_j} + \eta_{i_j}) - F(a_{i_j}) \le \sum_{i=1}^\infty F(b_i + \eta_i) - F(a_i).
\]
In the last step, we have simply added infinitely many nonnegative terms on the right side. We have used the assumption that $F$ is a non-decreasing function in this step.

We express all terms in terms of $\epsilon$ and simplify the inequality to
\begin{align*}
F(b) - F(a) - \epsilon &\le \sum_{i=1}^\infty \left( F(b_i) - F(a_i) + \frac{\epsilon}{2^i} \right) \\
F(b) - F(a) &\le 2\epsilon + \sum_{i=1}^\infty \big( F(b_i) - F(a_i) \big).
\end{align*}
Because $\epsilon$ could be arbitrarily small, we obtain
\[
F(b) - F(a) \le \sum_{i=1}^\infty \big( F(b_i) - F(a_i) \big).
\]
This proves that $\mu_0$ is $\sigma$-subadditive.

\textbf{Step 4.} $\mu_0$ is $\sigma$-finite.

We can partition the real number line as a union of $(m - 1, m]$, for $m \in \mathbb{Z}$. The pre-measure $\mu_0((m - 1, m])$ is finite for all $m$.

Combining all the previous steps, we see that $\mu_0$ is a $\sigma$-finite pre-measure on the algebra $\mathcal{B}_0$, which generates the Borel algebra on $\mathbb{R}$. By applying the measure extension theorem (Theorem 3.1), we obtain a unique measure on $\mathcal{B}(\mathbb{R})$ that extends $\mu_0$. \hfill $\blacksquare$